\section{Language and library choices}
\label{sec:languages}

The repository combines Rust and Python because annotation, machine-learning orchestration, and video I/O impose different constraints on responsiveness, dependency weight, and ecosystem access.

\subsection{Rust segmentator}

The desktop segmentator is implemented in Rust with \texttt{egui} and \texttt{eframe}. This stack was selected for a native polygon editor that watches dataset folders, refreshes image lists, and serialises YOLO segmentation labels beside each image with low UI latency. It corresponds to stage~3 in Figure~\ref{fig:pipeline14steps}. Dataset discovery, label parsing, and file watching are concentrated in \texttt{segmentator/src/dataset.rs}. Unit tests for label I/O live in inline \texttt{\#[cfg(test)]} modules next to the code they exercise.

\subsection{Python extractor and ML pipeline}

Two independent Python packages are managed with \texttt{uv}, each with its own \texttt{pyproject.toml} and lockfile. The extractor under \texttt{extractor/yolo\_raw\_extractor} targets Python~3.13 and handles frame extraction, segment cropping, and dataset augmentation (stages~1, 2, and~4 in Figure~\ref{fig:pipeline14steps}). The ML pipeline under \texttt{segmentator/ml\_pipeline} targets Python~3.11 for compatibility with Ultralytics and PyTorch and implements stages~5--14 of the same figure.

Training, validation, prediction, export, and video blur depend on the optional \texttt{[train]} extra (Ultralytics YOLO26-seg). ONNX export additionally requires the \texttt{[export]} extra. Layout and YAML utilities install after a plain \texttt{uv sync} without those extras, which keeps continuous integration lightweight. Video anonymisation reads and writes frames through OpenCV in \texttt{video.py}. Dataset hygiene, cleaning, iteration wrappers, and metric interpretation remain in standard-library-oriented modules with small, testable CLI entry points.

\subsection{Rationale for a hybrid layout}

Interactive labelling benefits from a compiled desktop application, whereas model training and batch CLIs benefit from the Python machine-learning ecosystem. The design replaces monolithic notebooks with path-based commands that can be composed in shell scripts or documented pipelines, which aligns with the staged layout shown in Figure~\ref{fig:pipeline14steps} and with reproducibility goals emphasised in the course material on clean code and pipelining.
